\documentclass[11pt,a4paper]{article}
\usepackage[utf8]{inputenc}
\usepackage[T2A]{fontenc}
\usepackage[russian]{babel}
\usepackage{amsmath, amssymb, amsfonts, amsthm}
\usepackage{geometry}
\usepackage{booktabs}
\usepackage{cite}
\usepackage{caption}
\usepackage{hyperref}
\usepackage{bm}

\geometry{top=2.0cm, bottom=2.0cm, left=2.0cm, right=2cm}
\title{\textbf{Азъ\\ Основания топологической эластодинамики вакуума: Волновая онтология микромира}}
\author{DM}
\date{\today}

\begin{document}
	
\section{Эластодинамика вакуума:\\ Тензор напряжений и спектр масс лептонов}

\subsection{Гиперупругий Лагранжиан и отказ от феноменологических потенциалов}
Классические топологические теории поля (такие как модель Скирма или Фаддеева-Ниеми) для предотвращения коллапса солитонов (обхода теоремы Деррика) вынуждены вводить феноменологические слагаемые высших порядков (члены Скирма). В рамках последовательной механики сплошных сред такой подход является искусственным. Стабильность макроскопического дефекта в упругом вакууме должна обеспечиваться естественным термодинамическим балансом — пределом текучести среды.

Для строгого описания мы переходим от абстрактного параметра порядка к векторному полю физических смещений континуума $\mathbf{u}(\mathbf{x}, t)$. Тензор упругих напряжений $T_{ij}$ в нелинейной среде (гиперупругости) математически строго разбивается на две ортогональные компоненты:
\begin{equation}
	T_{ij} = \mu_0 \delta_{ij} + S_{ij},
\end{equation}
где:
\begin{itemize}
	\item $\mu_0 = \frac{1}{3}\text{Tr}(T)$ — гидростатическое напряжение (эквивалент радиационного давления запертой волны, вызывающего объемное расширение/сжатие трубки).
	\item $S_{ij}$ — девиатор напряжений (тензор чистого сдвига), описывающий изменение формы (кручение, топологический спинорный твист) без изменения объема. По определению $\text{Tr}(S) = 0$.
\end{itemize}

Истинный фундаментальный Лагранжиан вакуумного солитона формулируется через инварианты этого тензора:
\begin{equation}
	\mathcal{L} = \frac{1}{2} \rho (\partial_t \mathbf{u})^2 - \left[ \frac{K}{2} I_1^2 + G J_2 \right] + \Lambda \left( 3\mu_0^2 - 2J_2 \right) + \mathcal{L}_{\text{gauge}}.
\end{equation}
Здесь $\frac{K}{2} I_1^2 + G J_2$ — классическая упругая энергия Гука (сжатие и сдвиг). 
Фундаментальным прорывом модели является множитель Лагранжа $\Lambda \left( 3\mu_0^2 - 2J_2 \right)$. Это слагаемое накладывает жесткую нелинейную связь на тензор напряжений, известную в механике сплошных сред как \textbf{Критерий текучести Губера — фон Мизеса}. 
Физический смысл этого критерия: локализованный фазовый вихрь (солитон) стабилен тогда и только тогда, когда упругая энергия изменения объема (давление волны) находится в строгом равновесии с упругой энергией изменения формы (топологическим кручением). При $2J_2 > 3\mu_0^2$ вакуум разрывается (сфалеронный пробой), при $2J_2 < 3\mu_0^2$ солитон рассеивается как безмассовое излучение.

\subsection{Тензорное происхождение 3 поколений и аналитический вывод формулы Коиде}
Поскольку масса частицы в гармоническом волновом приближении пропорциональна квадрату амплитуды локальной деформации упругой среды ($E \propto \lambda^2$), мы устанавливаем строгую связь: собственные значения тензора напряжений $\lambda_n$ являются квадратными корнями из наблюдаемых масс покоя лептонов ($\lambda_n = \sqrt{m_n}$).

Из линейной алгебры известно, что любой симметричный 3D-тензор ($3 \times 3$) обладает ровно тремя главными осями и тремя собственными значениями. Это математически детерминирует существование \textbf{строго трех поколений лептонов} как трех ортогональных резонансных мод (осей поляризации) единого тороидального дефекта вакуума. Никакого «четвертого поколения» не существует в силу трехмерности физического пространства тензора деформаций.

Собственные значения девиатора $s_n$ (где $\sqrt{m_n} = \mu_0 + s_n$) определяются характеристическим кубическим уравнением $\det(S_{ij} - s \mathbf{I}) = 0$. Так как $\text{Tr}(S)=0$, квадратичный член исчезает:
\begin{equation}
	s^3 - J_2 s - J_3 = 0.
\end{equation}
Аналитическое решение этого уравнения (формула Виета-Кардано) дает тригонометрический спектр собственных амплитуд:
\begin{equation}
	s_n = 2\sqrt{\frac{J_2}{3}} \cos\left( \delta + \frac{2\pi n}{3} \right), \quad n=0, 1, 2.
\end{equation}

Проверим, к какому спектру масс приводит стабильность солитона. Вычислим эмпирическую дробь Коиде $Q = \frac{\sum m_n}{(\sum \sqrt{m_n})^2}$ для нашего тензора:
\begin{itemize}
	\item Знаменатель: $\sum \sqrt{m_n} = \sum (\mu_0 + s_n) = 3\mu_0$ (так как $\sum s_n = 0$). Следовательно, $(\sum \sqrt{m_n})^2 = 9\mu_0^2$.
	\item Числитель (сумма масс): $\sum m_n = \sum (\mu_0 + s_n)^2 = 3\mu_0^2 + 2\mu_0\sum s_n + \sum s_n^2$. Поскольку $\sum s_n^2 \equiv 2J_2$ (по определению второго инварианта), сумма масс равна $3\mu_0^2 + 2J_2$.
\end{itemize}
Подставляем в формулу Коиде:
\begin{equation}
	Q = \frac{3\mu_0^2 + 2J_2}{9\mu_0^2} = \frac{1}{3} + \frac{2J_2}{9\mu_0^2}.
\end{equation}
Для стабильного солитона (согласно выведенному из Лагранжиана критерию Мизеса) вакуум жестко требует баланса $3\mu_0^2 = 2J_2$. Подставляя это условие в дробь, мы получаем точный аналитический результат:
\begin{equation}
	Q = \frac{1}{3} + \frac{3\mu_0^2}{9\mu_0^2} = \frac{1}{3} + \frac{1}{3} = \mathbf{\frac{2}{3}}.
\end{equation}
\textbf{Вывод:} Знаменитая формула масс Коиде ($Q=2/3$), связывающая массы электрона, мюона и тау-лептона, не является феноменологической нумерологией. Это точное уравнение состояния (термодинамический предел текучести) сплошной упругой среды в трехмерном пространстве. Лептоны — это три главных корня тензора напряжений стабильного вакуумного вихря.

\subsection{Универсальный предел сдвига $\alpha$ и макроскопическая геометрия $r/R$}
Матрица напряжений $S_{ij}$, генерирующая массы, зависит исключительно от локального угла наклона фазового тока (отношения кручения к давлению) $\gamma$.
В упругом вакууме стабильная стоячая волна может течь только под одним универсальным углом, при котором радиационное давление волны идеально уравновешивается магнитным/топологическим пинч-эффектом. Этот локальный предел текучести вакуума является абсолютной константой среды, геометрическим выражением которой служит постоянная тонкой структуры:
\begin{equation}
	\tan \gamma \equiv \alpha.
\end{equation}

Однако лептоны ($e, \mu, \tau$) характеризуются различной внутренней топологией — узлами $T(p,q)$ с индексами намотки фазы (в силу спинорной статистики фермионов $p_n=n, q_n=2n-1$).
Тангенс локального угла наклона спирали на торе задается отношением полоидального и тороидального путей волны:
\begin{equation}
	\tan \gamma = \frac{p \cdot (2\pi r)}{q \cdot (2\pi R)} = \frac{p}{q} \frac{r}{R}.
\end{equation}

Возникает фундаментальное противоречие: если бы пропорция $r/R$ была одинаковой для всех частиц, тензор деформаций зависел бы от $p$ и $q$, и лептоны не могли бы являться корнями единого кубического уравнения Коиде. 
Природа разрешает это противоречие механизмом \textbf{геометрической компенсации}. Вакуум жестко фиксирует локальный угол $\tan \gamma = \alpha$. Чтобы сохранить этот универсальный тензор при увеличении числа витков $q/p$, тор обязан макроскопически изменить толщину своего сечения:
\begin{equation}
	\frac{r_n}{R_n} = \frac{q_n}{p_n} \alpha = \frac{2n-1}{n} \alpha.
\end{equation}

Отсюда мы без феноменологических подгонок выводим точную геометрию всех трех лептонов:
\begin{itemize}
	\item \textbf{Электрон ($n=1, T(1,1)$):} $r/R = 1 \cdot \alpha$. Экстремально тонкое, "расслабленное" кольцо.
	\item \textbf{Мюон ($n=2, T(2,3)$):} $r/R = 1.5 \alpha$. Внутреннее кручение стягивает большой радиус, сечение распухает для сохранения угла текучести.
	\item \textbf{Тау-лептон ($n=3, T(3,5)$):} $r/R = 1.67 \alpha$. Максимально «толстый» тор. Дальнейшее усложнение узла ($n \ge 4$) потребовало бы такого распухания сечения, при котором исчезает центральное отверстие ($r \to R$), что топологически разрушает тор, делая 4-е поколение лептонов математически невозможным.
\end{itemize}

Различие масс ($m_e, m_\mu, m_\tau$) обусловлено исключительно этой геометрической подстройкой: локальный тензор напряжений для них АБСОЛЮТНО ИДЕНТИЧЕН (откуда вытекает формула Коиде), но макроскопические масштабы (комптоновские радиусы $R_n$) коллапсируют пропорционально росту энергии упругой фрустрации узла.

\section{Топологическая электродинамика: Заряд и Магнитный момент}

\subsection{Происхождение электрического заряда (Механизм Джулии-Зи)}
В рамках упругого вакуума электрический заряд не является фундаментальной «сущностью», помещенной на частицу. Заряд — это топологический поток (циркуляция) градиента фазы $\Phi$ внутреннего спинорного тока, сшиваемый с односвязным внешним 3D-пространством.

Вектор скорости фазовой волны на торическом узле имеет две компоненты: тороидальную (вдоль большого кольца) и полоидальную (поперек сечения трубки). 
Тороидальный ток замкнут сам на себя внутри 3D-пространства. Он не излучает радиального поля смещения наружу. Напротив, полоидальный ток (выворачивающий тор наизнанку) топологически не может быть замкнут без создания радиального градиента напряжений в окружающем эфире. 
Внешний вакуум воспринимает этот полоидальный фазовый поток как пульсирующий монополь (кулоновское поле). Электрический заряд строго квантуется через Индекс Зацепления (Linking Number) $L_k$ полоидального контура с любой замкнутой 2D-сферой, охватывающей частицу:
\begin{equation}
	Q_{electric} \propto \oint_{poloidal} \nabla \Phi \cdot d\mathbf{l} \implies L_k = 1.
\end{equation}
Поскольку все три поколения лептонов $(e, \mu, \tau)$ являются топологически гомологичными торами (их полоидальные витки, независимо от числа $p$, имеют строго одинаковый индекс зацепления $1$ с внешней оболочкой), их измеряемый электрический заряд абсолютно идентичен и равен $-1e$.

\subsection{Строгий вывод Магнитного Момента}
Магнитный момент $\mu_z$ генерируется тороидальной (продольной) компонентой фазового тока. Рассчитаем его классически: $\mu_z = I_{tor} \cdot A_{eff}$.

\textbf{1. Эффективная площадь $A_{eff}$:}
Волна распределена по поверхности тора $T(R, r)$. Интегрирование контура по проекции на плоскость $XY$ дает строгую аналитическую площадь:
\begin{equation}
	A_{eff} = \frac{1}{2\pi} \int_0^{2\pi} \pi (R^2 + 2Rr\cos\theta + r^2\cos^2\theta) d\theta = \pi R^2 \left( 1 + \frac{1}{2} \left(\frac{r}{R}\right)^2 \right).
\end{equation}

\textbf{2. Эффективный ток $I_{tor}$:}
Волна движется по поверхности тора со скоростью света $c$. Вектор скорости направлен под углом кручения $\gamma$ к оси трубки (напомним, предел текучести вакуума требует $\tan \gamma \equiv \alpha$). 
Скорость волны вдоль тороидального направления (оси кольца) равна:
\begin{equation}
	v_{tor} = c \cos \gamma = \frac{c}{\sqrt{1 + \alpha^2}}.
\end{equation}
Время, за которое фронт волны совершает один полный оборот вокруг оси $Z$ (период тороидального обхода), равно $T_{tor} = 2\pi R / v_{tor}$. Средний макроскопический ток, формирующий магнитный диполь, равен заряду $e$, деленному на этот период:
\begin{equation}
	I_{tor} = \frac{e v_{tor}}{2\pi R} = \frac{e c \cos \gamma}{2\pi R}.
\end{equation}
*(Примечание: Топологическое число витков $q$ здесь строго сокращается, так как волна проходит через сечение $q$ раз за один полный цикл узла, делая средний ток зависящим только от базового тороидального периода).*

\textbf{3. Комптоновский резонанс и масса:}
Макроскопический радиус тора $R$ определяется условием стоячей волны (балансом энергии). Для образования стабильного солитона радиус обязан резонировать с массой частицы (соответствовать комптоновскому радиусу переносчика):
\begin{equation}
	R = \frac{\hbar}{mc}.
\end{equation}

\textbf{4. Формула голого магнитного момента:}
Перемножаем ток и площадь, подставляя $R = \hbar/mc$:
\begin{equation}
	\mu_z = \left( \frac{e c \cos \gamma}{2\pi R} \right) \times \left[ \pi R^2 \left( 1 + \frac{1}{2} \left(\frac{r}{R}\right)^2 \right) \right] = \frac{e c R \cos \gamma}{2} \left( 1 + \frac{1}{2} \left(\frac{r}{R}\right)^2 \right).
\end{equation}
\begin{equation}
	\mu_z = \left( \frac{e \hbar}{2m} \right) \cos \gamma \left( 1 + \frac{1}{2} \left(\frac{r}{R}\right)^2 \right) = \mu_B \cos \gamma \left( 1 + \frac{1}{2} \left(\frac{r}{R}\right)^2 \right).
\end{equation}

\subsection{$g=2$ для Электрона}
Аналитически проверим эту геометрическую формулу для базового состояния вакуума — электрона.
Для электрона ($n=1$) мы ранее строго вывели макроскопическую пропорцию: $r/R = \alpha$.
Подставляем $\cos \gamma = (1 + \alpha^2)^{-1/2}$ и $r/R = \alpha$ в уравнение магнитного момента:
\begin{equation}
	\frac{\mu_e}{\mu_B} = \frac{1 + \frac{1}{2} \alpha^2}{\sqrt{1 + \alpha^2}}.
\end{equation}
Разложим это выражение в ряд Тейлора по степеням малой величины $\alpha$:
\begin{equation}
	\frac{1}{\sqrt{1 + \alpha^2}} \approx 1 - \frac{1}{2}\alpha^2 + \frac{3}{8}\alpha^4 - \dots
\end{equation}
Умножаем на площадь сечения $(1 + \frac{1}{2}\alpha^2)$:
\begin{equation}
	\frac{\mu_e}{\mu_B} = \left( 1 + \frac{1}{2}\alpha^2 \right) \left( 1 - \frac{1}{2}\alpha^2 + \frac{3}{8}\alpha^4 \right) = 1 - \frac{1}{2}\alpha^2 + \frac{1}{2}\alpha^2 - \frac{1}{4}\alpha^4 + \frac{3}{8}\alpha^4 = 1 + \mathbf{\frac{1}{8}\alpha^4}.
\end{equation}

\textbf{Фундаментальный результат:} Члены с $\alpha^2$ \textbf{АБСОЛЮТНО СОКРАТИЛИСЬ}! 
Уменьшение скорости тороидального тока (за счет скручивания $\cos \gamma$) ИДЕАЛЬНО скомпенсировано увеличением эффективной площади из-за толщины тора ($r^2$). 
Геометрический магнитный момент голого электрона (без учета радиационного трения вакуума) равен дираковскому $g=2$ с фантастической точностью (ошибка возникает только в 4-м порядке $\alpha^4$, что составляет $\sim 3 \cdot 10^{-10}$). 
Это доказывает, что дираковский $g$-фактор не требует постулирования "точечной частицы", а является точным классическим свойством вихревого кольца $T(1,1)$ на пределе текучести среды $\tan\gamma=\alpha$.

\subsection{Геометрическая природа массозависимых аномалий ($\mu, \tau$)}
В отличие от электрона, тяжелые лептоны вынуждены "распухать", чтобы сохранить локальный угол $\alpha$. Для $n$-го поколения мы вывели $r_n / R_n = k_n \alpha$, где $k_n = \frac{2n-1}{n}$.
Магнитный момент для них (в ряду Тейлора):
\begin{equation}
	\frac{\mu_n}{\mu_B} \approx \left( 1 + \frac{1}{2} k_n^2 \alpha^2 \right) \left( 1 - \frac{1}{2}\alpha^2 \right) = 1 + \frac{1}{2}(k_n^2 - 1)\alpha^2.
\end{equation}
Поскольку для мюона ($k_2 = 1.5$) и тау-лептона ($k_3 = 1.67$) топологический фактор $k_n > 1$, они получают \textbf{строго положительную добавку} к своему базовому магнитному моменту ($\sim +0.625\alpha^2$ для мюона). 
В Стандартной Модели (КЭД) эта массозависимая разница высчитывается через виртуальные петли частиц разных поколений. В топологической модели мы получаем математически строгий эквивалент: мюон имеет более высокий голый магнитный момент просто потому, что его топологический узел $T(2,3)$ требует более толстого сечения тора, что квадратично увеличивает макроскопическую площадь токового кольца.

\end{document}